\section{Parameters and calibration}
\label{sec:params}

\subsection{Provenance of the inputs}

No measurements were made for this study. Every input therefore comes from one of five sources, and Table~\ref{tab:params} states which one applies. A value is labelled L when it was reported for the glass-ceramic of the IBM process or for copper itself; A when it was taken from an analogous material, usually a commercial copper-compatible LTCC; C when it was calibrated here against published process observations (Section~\ref{sec:calibration}); D when it follows from other entries of the table; and E when it is an engineering estimate. Most estimated and analogue values that matter enter the sensitivity analysis of Section~\ref{sec:res_sens}, with the ranges given in the table; three that do not were checked one at a time.

The reference glass is the MgO-rich cordierite-type glass used by IBM for its copper-compatible substrates (glass \#12: 52.5 SiO$_2$, 22.0 MgO, 22.0 Al$_2$O$_3$, 1.5 P$_2$O$_5$, 0.5 B$_2$O$_3$, 1.5 ZrO$_2$, wt\%) \cite{herron1980method,kumar1981glassceramic}. It was chosen because, among copper-compatible glass-ceramics, it is the one for which the firing process, the atmosphere and the reasons behind the chosen temperatures are best documented. We could not obtain published viscosity data for it, so its viscosity was calibrated.

\begin{table*}[!tp]
\centering
\caption{Model inputs. Basis: L, literature value for this material; A, analogue material; C, calibrated in Section~\ref{sec:calibration}; D, derived; E, estimate. Ranges in the last column are those used in the Morris screening.}
\label{tab:params}
\footnotesize
\begin{tabular}{@{}llllp{7.2cm}@{}}
\toprule
Quantity & Symbol & Value & Basis & Source, remark or screening range \\
\midrule
\multicolumn{5}{@{}l}{\textit{Glass-ceramic body}}\\
Median particle size & $d_{50}$ & 3.0~\textmu m & A & IBM glass powders 2--7~\textmu m \cite{kumar1981glassceramic}; 2--7~\textmu m \\
Size span $(d_{90}-d_{10})/d_{50}$ & -- & 1.5 & E & \\
Solids loading of the slurry & $\phi_{\mathrm{GC}}$ & 0.45 & E & 0.40--0.50 \\
Density of the dense body & $\rho_{\mathrm{s}}$ & 2600~kg\,m$^{-3}$ & A & dense $\alpha$-cordierite glass-ceramic, 2620~kg\,m$^{-3}$ \cite{yu2022synthesis} \\
Surface energy & $\gamma$ & 0.36~J\,m$^{-2}$ & L & calculated for a cordierite-type glass \cite{giess1984isothermal}; 0.24--0.36, since water vapour lowers the surface tension of silicate glass \cite{parikh1958effect} \\
Glass transition ($10^{12}$~Pa\,s) & $T_{\mathrm{g}}$ & 734~$^\circ$C & C & 714--754~$^\circ$C \\
Fragility index & $m$ & 32.4 & C & 28--40 \\
High-temperature viscosity limit & $\log_{10}\eta_\infty$ & $-3$ & L & \cite{mauro2009viscosity} \\
Crystallisation activation energy & $E_{\mathrm{c}}$ & 303~kJ\,mol$^{-1}$ & L & IBM-type high-cordierite glass \cite{watanabe1994crystallization}; 330--440~kJ\,mol$^{-1}$ in other MgO--Al$_2$O$_3$--SiO$_2$ glasses \cite{sun2022crystallization,wang2026effect}; 250--470~kJ\,mol$^{-1}$ \\
Avrami exponent & $n$ & 2 & E & 1.5--3.7; 2.8--3.7 in a related MgO--Al$_2$O$_3$--SiO$_2$ glass \cite{wang2026effect} \\
Exotherm peak at 10~K\,min$^{-1}$ & $T_{\mathrm{p,c}}$ & 1040~$^\circ$C & C & 1000--1080~$^\circ$C \\
Maximum crystal packing & $X_{\mathrm{m}}$ & 0.64 & E & random close packing \\
Density at pore closure & $\rho_{\mathrm{cl}}$ & 0.92 & E & checked at 0.88 and 0.95 \\
Thermal conductivity, dense & $k$ & 2.8~W\,m$^{-1}$\,K$^{-1}$ & A & copper co-fired LTCC GL570 \cite{kyocera2026material} \\
Young's modulus, dense & $E_{\mathrm{GC}}$ & 128~GPa & A & GL570 \cite{kyocera2026material} \\
Poisson's ratio & $\nu_{\mathrm{GC}}$ & 0.25 & E & \\
Thermal expansion coefficient & $\alpha_{\mathrm{GC}}$ & 3.0~ppm\,K$^{-1}$ & A & production glass-ceramic substrate \cite{knickerbocker2002advanced}; 1.8~ppm\,K$^{-1}$ for glass \#12 fired in H$_2$/H$_2$O \cite{herron1980method} \\
Flexural strength & $\sigma_{\mathrm{s,GC}}$ & 210~MPa & L & glass \#12 fired in H$_2$/H$_2$O \cite{herron1980method} \\
Stress-free temperature & $T_{\mathrm{sf}}$ & 690~$^\circ$C & D & strain point ($10^{13.5}$~Pa\,s) of the calibrated glass \\
\midrule
\multicolumn{5}{@{}l}{\textit{Copper conductor}}\\
Median particle size & $d_{50}$ & 1--20~\textmu m & -- & design variable; baseline 3~\textmu m \\
Size span & -- & 1.2 & E & \\
Solids loading & $\phi_{\mathrm{Cu}}$ & 0.40--0.60 & -- & design variable; baseline 0.50 \\
Inert filler fraction & $f$ & 0--0.40 & -- & design variable; baseline 0 \\
Viscosity factor & $f_\eta$ & 0.28 & C & \ref{app:copper}, from \cite{roumanie2021influence}; 0.1--1.0 \\
Lattice and boundary diffusion & $D_{\mathrm{v}}$, $\delta D_{\mathrm{b}}$ & -- & L & \cite{frost1982deformation} \\
Carbon inhibition scale & $C_{\mathrm{inh}}$ & 200~ppm & E & checked at 50 and 1000~ppm \\
Young's modulus & $E_{\mathrm{Cu}}$ & 117~GPa & L & oxygen-free copper \cite{aurubisc10200} \\
Poisson's ratio & $\nu_{\mathrm{Cu}}$ & 0.34 & E & \\
Thermal expansion coefficient & $\alpha_{\mathrm{Cu}}(T)$ & 16.5--23.7~ppm\,K$^{-1}$ & L & instantaneous values, 293--1100~K \cite{kayelaby235} \\
Yield stress & $\sigma_{\mathrm{y}}(T)$ & 69~MPa at 20~$^\circ$C & L & annealed temper \cite{aurubisc10200}; softening $1-T^{*1.09}$ \cite{johnson1985fracture,banerjee2005evaluation} \\
\midrule
\multicolumn{5}{@{}l}{\textit{Binder, char and atmosphere}}\\
Char yield of network and backbone & $\chi$ & 0.05 & E & 0.02--0.10 \\
Gasification half-life, 800~$^\circ$C, reference gas & $t_{1/2}$ & 2~h & E & 0.5--8~h \\
Gasification activation energy & $E_{\mathrm{g}}$ & 200~kJ\,mol$^{-1}$ & A & order of steam gasification of carbons \cite{walker1959gas}; 160--240~kJ\,mol$^{-1}$ \\
Hydrogen inhibition constant & $K_{\mathrm{H}}$ & 25 & E & form after \cite{barrio2001steam}; checked at 5 and 100 \\
Cu/Cu$_2$O and C/H$_2$O equilibria & $K_{\mathrm{Cu}}$, $K_{\mathrm{st}}$ & -- & L & NIST-JANAF \cite{chase1998janaf} \\
Safety factor on the Cu/Cu$_2$O limit & $s$ & 0.2 & E & \\
\midrule
\multicolumn{5}{@{}l}{\textit{Geometry}}\\
Slab half-thickness & -- & 2.0 / 0.15~mm & E & glass-ceramic / copper \\
Reference bilayer & $h_{\mathrm{Cu}}$, $h_{\mathrm{GC}}$ & 0.1, 1.0~mm & E & 20~mm span \\
\bottomrule
\end{tabular}
\end{table*}

\subsection{Calibration of the glass against the IBM process}
\label{sec:calibration}

The patents that describe the IBM steam co-sintering process state, for glass \#12, why the burnout hold sits where it does. The hold is placed at $785 \pm 10$~$^\circ$C, between the annealing and softening points of the glass. At 750~$^\circ$C carbon removal would take ``prohibitively excessive amounts of time''; at 830~$^\circ$C the glass pores close and trap water and binder residue \cite{herron1980method}. A later IBM patent on the co-sintering of multilayer ceramic substrates states that 33--50\% of the total fired shrinkage occurs during the burnout step \cite{flaitz1992method}; we assume that this holds for glass \#12, and a third places the onset of crystallisation at about 900~$^\circ$C, after the burnout hold and before the crystallisation hold at about 960~$^\circ$C \cite{dubetsky1982process}. We simulated the full worked schedule of Ref.~\cite{herron1980method} (N$_2$ to 200~$^\circ$C, H$_2$/H$_2$O to 450~$^\circ$C, 2.9~K\,min$^{-1}$ to the hold, 6~h at H$_2$/H$_2$O $= 10^{-4}$, 0.5~h in N$_2$, 2.1~K\,min$^{-1}$ to 960~$^\circ$C, 2~h) and used these statements as targets. The patent specifies the H$_2$/H$_2$O ratio but not the dilution; we assumed 50\% steam in nitrogen.

Two of them fix the viscosity. Requiring that a 6~h hold at 780~$^\circ$C produce 40\% of the fired shrinkage with the pores still open, and that the pores close within 2~h at 830~$^\circ$C, gives $\log_{10}\eta = 10.62$ at 780~$^\circ$C and 9.37 at 830~$^\circ$C. The MYEGA curve through these two points, with $\log_{10}\eta_\infty = -3$, has $T_{\mathrm{g}} = 733.6$~$^\circ$C and $m = 32.4$. The remaining observations were not used in the fit and serve as checks. Between 800 and 860~$^\circ$C the calibrated viscosity falls by one decade per 43~K. Giess et al.\ found that the sintering of an IBM cordierite-type glass over the same range followed the viscosity of the glass, which fell by about one decade per 40~K \cite{giess1984isothermal}. The annealing point, taken as $10^{12}$~Pa\,s, lies at 734~$^\circ$C and the Littleton softening point ($10^{6.6}$~Pa\,s) at 980~$^\circ$C, so the 785~$^\circ$C hold lies between them as the patent requires; so does the dilatometric softening point, near $10^{10}$~Pa\,s, which the calibrated glass reaches at about 803~$^\circ$C. Stoichiometric cordierite glass is more viscous by 3.7 decades at 780~$^\circ$C, 1.0 at 900~$^\circ$C and 0.4 at 960~$^\circ$C \cite{reinsch2008crystal}. Glass \#12 contains more MgO and less Al$_2$O$_3$ than cordierite, together with B$_2$O$_3$ and P$_2$O$_5$, so a lower viscosity is plausible, but this comparison is not a test of the fit.

The crystallisation kinetics were treated in the same way. With the activation energy measured on an IBM-type glass ($E_{\mathrm{c}} = 303$~kJ\,mol$^{-1}$ \cite{watanabe1994crystallization}) and $n = 2$, placing the onset of crystallisation (5\% crystallised) at 900~$^\circ$C in the worked schedule requires an exotherm peak at 1040~$^\circ$C at 10~K\,min$^{-1}$. The simulated exotherm then spans 978--1066~$^\circ$C, a width of about 90~K, comparable to the 80--100~$^\circ$C span of sintering temperatures that the IBM patent finds covered by the crystallisation exotherm of its glasses \cite{kumar1981glassceramic}. Differential thermal analysis of other MgO--Al$_2$O$_3$--SiO$_2$ glasses places the $\alpha$-cordierite peak at 1006--1032~$^\circ$C \cite{yu2022synthesis,wang2026effect}.

\begin{table}[!tp]
\centering
\caption{The IBM worked schedule \cite{herron1980method} simulated at three hold temperatures. Shrinkage is the fraction of the total fired shrinkage reached by the end of the 6~h hold; carbon is the residual carbon in the glass-ceramic at the end of the hold.}
\label{tab:ibm}
\footnotesize
\begin{tabular}{@{}lccc@{}}
\toprule
Hold temperature & 750~$^\circ$C & 785~$^\circ$C & 830~$^\circ$C \\
\midrule
Shrinkage during hold & 7\% & 47\% & 100\% \\
Pores at end of hold & open & open & closed \\
Carbon at end of hold (ppm) & 33 & 0.006 & 0.07 \\
Final relative density & 0.990 & 0.990 & 0.990 \\
Final crystallised fraction & 1.00 & 1.00 & 1.00 \\
\midrule
Statement in \cite{herron1980method} & too slow & chosen & pores close \\
\bottomrule
\end{tabular}
\end{table}

Table~\ref{tab:ibm} summarises the result. The 785~$^\circ$C hold removes the char and leaves 47\% of the shrinkage to the burnout step, inside the 33--50\% reported for the process, while the pores stay open. At 830~$^\circ$C the body densifies completely during the hold and seals its pores. These two rows reproduce the calibration targets and are therefore not independent evidence. The 750~$^\circ$C row is independent: it depends on the gasification kinetics, which were not fitted, and it leaves four orders of magnitude more carbon than the 785~$^\circ$C hold. The model does not reproduce the absolute time scale of the patent's statement. With the estimated gasification rate, 33~ppm of carbon remain after 6~h at 750~$^\circ$C, and whether that would have been acceptable cannot be judged, because no carbon specification was published. Nor does the model reproduce the reason the patent gives for avoiding 830~$^\circ$C: the carbon has gone before the pores close, and the trapped water that the patent mentions is not modelled. For all three holds the final body is 99\% dense and fully crystallised, and its free linear shrinkage is 23.1\%.
